% !TEX root = ../BeamerTemplate.tex
%%%%%%%%%%%%%%%%%%%%%%%%%%%%%%%%%%%%%%%%%%%%%%%%%%%%%%%%%%%%%%%%%%%%%%%%%%%%%%%%%%

\section{Conclusions and Future Work}

\begin{frame}[t]{The Four Research Questions Receive Conditional Answers}
\scriptsize
\begin{columns}[T,onlytextwidth]
  \column{0.49\textwidth}
  \begin{block}{RQ1 --- Physical realization: supported}
    The $20\times20$ transmitarray produces a realized focus near 14~mm, a
    focal arc, and angle-dependent far-field responses.
  \end{block}
  \begin{block}{RQ3 --- Mobile comparison: no universal gain}
    Under equal weighting, the baseline has higher median link and constructed
    MIMO capacity; the lens has stronger typical-block decorrelation.
  \end{block}

  \column{0.49\textwidth}
  \begin{block}{RQ2 --- Cross-solver trend: mostly supported}
    All solvers predict matched-link reinforcement, cross-link suppression,
    lower correlation, and higher raw capacity; normalized gains differ.
  \end{block}
  \begin{block}{RQ4 --- Higher order: gain with lower efficiency}
    Capacity grows from $3\times3$ to $9\times9$, while conditioning, rank
    fraction, and capacity per dimension deteriorate.
  \end{block}
\end{columns}
\Takeaway{The spatial transformation is real; its system benefit depends on
receiver operation and geometry.}
\SourceNote{Chapter 4: Principal Findings}
\end{frame}


\begin{frame}[t]{Principal Contributions}
\scriptsize
\begin{columns}[T,onlytextwidth]
  \column{0.49\textwidth}
  \begin{block}{Physical and computational bridge}
    \raggedright
    \begin{itemize}
      \setlength{\itemsep}{0.30em}
      \item Full-wave characterized planar $20\times20$ transmitarray and
            focal-plane RX arrangement at 38~GHz.
      \item Common raw/normalized MIMO framework for HFSS S-parameters and
            Sionna-RT channel responses.
      \item Three-solver evidence with 92.6\% agreement on the sign of the
            link-level lens effect.
    \end{itemize}
  \end{block}

  \column{0.49\textwidth}
  \begin{block}{System-level evidence}
    \raggedright
    \begin{itemize}
      \setlength{\itemsep}{0.30em}
      \item Paired mobile study explaining why local decorrelation does not
            automatically increase aggregate capacity.
      \item Explicit comparison of global, waypoint-local, and block-median
            decorrelation estimators.
      \item Higher-order evidence of absolute capacity growth with declining
            per-dimension efficiency.
    \end{itemize}
  \end{block}
\end{columns}
\SourceNote{Chapter 4: Contributions}
\end{frame}


\begin{frame}[t]{Overall Conclusion: Focusing Is Necessary but Not Sufficient}
\small
\begin{columns}[T,onlytextwidth]
  \column{0.31\textwidth}
  \begin{block}{What works}
    Passive focusing and angular filtering reshape channel amplitudes and
    improve local separability.
  \end{block}
  \column{0.31\textwidth}
  \begin{block}{What limits it}
    Misalignment, equal branch weighting, and geometry-limited singular modes
    can erase the link-level benefit.
  \end{block}
  \column{0.31\textwidth}
  \begin{block}{What is required}
    Placement, channel geometry, normalization, and beam management must be
    designed jointly.
  \end{block}
\end{columns}

\vspace{0.35cm}
\begin{beamercolorbox}[rounded=true,sep=1.0ex,wd=\linewidth]{takeaway}
  \centering\large\bfseries
  The transmitarray is an effective passive spatial transformer --- not an
  automatic MIMO-capacity multiplier.
\end{beamercolorbox}
\SourceNote{Chapter 4: Overall Conclusion}
\end{frame}


\begin{frame}[t]{Limitations I: Physical Validation and Channel Construction}
\scriptsize
\begin{columns}[T,onlytextwidth]
  \column{0.49\textwidth}
  \begin{block}{Physical and numerical validation}
    \begin{itemize}
      \setlength{\itemsep}{0.15em}
      \item Simulation only; no fabricated prototype or coherent OTA channel.
      \item Solver formulations and absolute reference quantities differ.
      \item Realized 14~mm focus differs from the 30~mm design target.
    \end{itemize}
  \end{block}

  \column{0.49\textwidth}
  \begin{block}{Receiver construction}
    \begin{itemize}
      \setlength{\itemsep}{0.15em}
      \item $N\times N$ matrices stack sequential $N\times1$ simulations.
      \item Mutual coupling, common phase, switching overhead, and joint RX
            noise are omitted.
      \item The reconstructed matrix is a controlled proxy, not a simultaneous
            measured multiport channel.
    \end{itemize}
  \end{block}
\end{columns}
\CautionNote{Absolute values from different solvers and channel constructions
are not directly interchangeable.}
\SourceNote{Chapter 4: Limitations}
\end{frame}


\begin{frame}[t]{Limitations II: Fair Comparison and Reproducibility}
\scriptsize
\begin{columns}[T,onlytextwidth]
  \column{0.49\textwidth}
  \begin{block}{Comparison confounds}
    \begin{itemize}
      \setlength{\itemsep}{0.15em}
      \item Lens and baseline designs differ in both placement and pattern.
      \item Order study changes power, aperture, RX count, and angular grid.
      \item Only the $0^\circ$ receive branch is common across all array orders.
    \end{itemize}
  \end{block}

  \column{0.49\textwidth}
  \begin{block}{Sampling and reproducibility}
    \begin{itemize}
      \setlength{\itemsep}{0.15em}
      \item One room and controlled trajectories; no multi-seed confidence
            intervals or convergence sweep.
      \item 100{,}000 ray samples/source; outage is saturated.
      \item Source notebooks and complete complex tensors are absent from the
            current trajectory archive.
    \end{itemize}
  \end{block}
\end{columns}
\CautionNote{Evidence supports controlled mechanisms, not universal deployment
performance.}
\SourceNote{Chapter 4: Limitations}
\end{frame}


\begin{frame}[t]{Future Work: Turn Angular Selectivity into Operational Gain}
\scriptsize
\begin{columns}[T,onlytextwidth]
  \column{0.49\textwidth}
  \begin{block}{1. Measure the physical system}
    Fabricate the feed and transmitarray; measure focal distance, patterns, and
    simultaneous coherent multiport channels.
  \end{block}
  \begin{block}{3. Make comparisons causal and fair}
    Isolate placement from pattern, then hold total TX power, aperture, and RX
    angle set fixed across orders.
  \end{block}

  \column{0.49\textwidth}
  \begin{block}{2. Add adaptive beam management}
    Evaluate selection, switching, combining, and hybrid beamforming with
    training overhead and alignment error.
  \end{block}
  \begin{block}{4. Broaden and reproduce the evidence}
    Retain tensors and notebooks; add seeds, convergence tests, blockage,
    mobility, richer scenes, multi-user precoding, and orders $N>9$.
  \end{block}
\end{columns}

\vspace{0.10cm}
\begin{beamercolorbox}[rounded=true,sep=0.75ex,wd=\linewidth]{takeaway}
  \centering
  \textbf{Joint target:}\quad lens design $+$ channel geometry $+$ RX placement
  $+$ beam-management algorithm.
\end{beamercolorbox}
\SourceNote{Chapter 4: Future Work}
\end{frame}
